\documentclass{ltjsbook}
\usepackage{docmute} 
\usepackage{../../myStyle} 

\begin{document}
\setcounter{chapter}{7}
\chapter{システム論からみた社会心理学}
\label{sp08_system}

\section{はじめに}
前回までで、先週は学際休みだったんですけど、社会心理学の現代に至るまでの歴史的なところはお話ししたつもりなんですが、今日からは話題が変わります。2つ目の話題に行こうと思っています。最初の授業の時にも一応3部作でという話はしてましたが、その2部目に今日から突入します。今日のお話はシステム論というところです。

このシステム論の話も多分2回、いや、2回は無理かな、3、4回はかかるかなと思っています。システムって，日本語で「系」という字を書くんですね。このシステム論というのは、どこかの学部で、またはどこかの学科で体系的に教えられるということは，おそらくほとんどないんじゃないかなと思います。唯一、僕は昔看護学校で非常勤したことがあるんですが，そこにはシステム論の授業がありました。その時もシステム論というのは本当にシステム論ですかと、それだけで単体の授業として成立するんですかと聞いたぐらい、ここだけ取り出して授業にしているところは少ないんじゃないかなと思います。

一体システム論というのは何の話をするのかというと、科学の具体的な研究テーマやリサーチクエスチョンがある話というよりは、科学観、ものの見方に基づいて論じられるものです。システム論としての共通した問題意識というのはあるんですけれども、それが何学なのかというのがわかりません。どこに分類して良いかがわからないというべきかな。

この授業は社会心理学の授業ですが、社会心理学の枠組みの中でシステム論を中心に扱うことは多分少ないんじゃないかなと思います。

今日はこのシステム論の歴史の話をします。システム論の歴史というのは、スタートは生物学、その後化学の方で発展し、その次に神経科学の方で発展しました。その後コンピュータ科学、社会心理学で発展した，と私は考えています。つまり，最終的には人文社会科学の枠組みでこれを話そうと思っています。ことほど左様に、学問領域を移動しながら学際的なものの見方に関するもので、見立ての学という側面があるものですから、ちょっと正体がつかみにくくて、教えた方もこれですよと、これがわかったらこういうことができるというのが伝えにくい形になっています。

歴史の話は社会心理学の歴史を追いかけてきたから良かったんですが、今回はそういったものの見方になりますので、ますます抽象的な変な話になるかなと思います。なぜ冒頭にこのようなことを言うかというと、いきなりシステム論を言い出して、しゃべり出してもいいんですが、ある程度のオリエンテーションがないと、受講生の皆さんも何を聞いてるんだろうとなるんじゃないかなと思ったからなんですね。

さて，これから何回かに分けてシステム論のお話をしますが、確実にお約束できるのは、私は社会心理学を考える上で、このシステム論の考え方が必要だと思っているという点です。このシステム論的な観点から社会心理学を考えるということをしたい、皆さんにもしてもらいたいんです。なので、ちょっと変な、例えば生物学だとか化学の話をしているなと思ったとしても、そこの根底にあるのは社会心理学を捉えるためのものの見方ということです。「システムとして社会心理学を考える」という軸はブレずに行きたいと思っておりますので、よろしくお付き合いいただければと思います。

\section{システム論の背景}
さて、システム論なんですが、これはどういった背景から出てきているかというと、単純に言って、「生きてるって何だろう」ということなんです。「生命」とは何か、ということを考えるための科学観、学問観としてスタートしています。生命とは何か、生きてるって何か、命とは何か、心とは何かというのは，非常に大きな問いです。しかし、これを考えるときに、「神様がくれたものだ」という考え方ではなく、科学的に考えるためにはどういった観点が必要なのか、というところからお話をしていきたいと思います。

さて、この授業の2回目か3回目くらいに、現代科学の基盤ということで、4つほど特徴を挙げさせていただきました。1つ目が実験的なアプローチです。ルネサンス前の時代に、ケプラーやガリレオなどは、社会で実際の観測データに基づいて、聖書に書いてあることを全面的に信じるのではなく、ピサの斜塔から鉄球を落とすといった実験をすることによって、本当に書かれていることが確かなのか確認しようというアプローチをとっていました。それまでは宗教的、キリスト中心的、教会中心的な考え方だったのですが、そこから自然科学的、実験的なアプローチが出てきたというのが、現代科学の一つの特徴だとお話しました。

これはこれでいいのですが、残る3つの特徴がありました。それは何かというと、1つは数学的還元主義、数値主義といった考え方です。ニュートンなどに代表される考え方ですが、それが一体研究しようとしている対象、それが一体何なのかということではなく、それをどのようにして表現するか、どのように表すかという方法を数学的な関数で表現しましょうという考え方です。特にここを数学的に表現できるということです。数式で表現しましょうというのが、近代科学の一つのやり方です。数式というのは、文系の人はあまり数学的な表現を好みませんが、それも言語の一種で、非常に厳密性の高い言語なので、数学で表現するというのがサイエンスの知識の積み重ねには必要だったわけです。

それと、3つ目として、機械論的な発想というのを挙げておきました。デカルトが提案した確率、というのは考え方の観点の確率なのですが、方法論の中に書いてあったように、いろんな物事の考え方を提案します。近代的な考え方を提案します。デカルトが言いたい機械論的な発想というのは、こういう原因があったらこういう結果にたどり着きますよね、こういうメカニズムで、この順番で物事が動きますねという因果的なプロセスを考えるという発想です。

何だかわからないけど、急に突然何かが現れるんだというのではなくて、おそらく原因があるはずです。現象には原因があるだろうと。何かのメカニズムがあるだろうと。マクスウェルのデーモンなんかという例え話がありますが、マクスウェルが考えたのは、ある特定の時間の空間にある全ての分子、原子の位置と速度、運動が分かる悪魔がいたら、この先何が起こるかを完全に予測できるだろうという考え方です。実際にそういう悪魔がいるかどうかという話は
抜きにして、その発想です。すべての条件を確実に全部把握すれば、因果のつながりで未来までしっかりと寸分違わず予測できますよ。なぜなら、それは物が動くには不思議な神秘的な力がかかっているのではなく、メカニズムがあるからです。そういうメカニズムに基づく考え方、というのがこの機械論的な発想です。これが近代科学の特徴の三つ目。そして四つ目は、これもデカルトの確立ですが、分解と統合、分立統合でしたか、大きな問題は小さなパートに分けて、それを組み立てていけば、最終的には全体を理解したことになる、という考え方ですね。これは要素還元主義という風に言われます。小さな要素に還元して、大きな複雑な問題は小分けにして考えていきましょう、というスタイルですね。

ということで、今挙げました実験的なアプローチ、数値主義、数学的な現象主義、機械論的な発想、因果的な発想と要素還元主義、これが近代科学の発想、考え方の根本なんですけれども、このアプローチでもっても捉えられないのが、生命というものです。生きているとはどういうことか、ということを考えるときに、今言ったようなアプローチが効かないわけですね。このアプローチでは生命ということの解明ができないというのが大きな問題なわけです。

何でもわかると、すぐにお分かりいただけるかと思うんですが、要素還元主義的に生命にアプローチしようとして、生きている人間をズバーッと切って、「さあどこが生きている要素なのかな」と言って、例えば上半身とか下半身に切り分けてみました。「どっちが生きているんだろう」と思ったら、両方死んでしまいました、みたいなことがあるわけですね。これは当たり前です。生きているっていうのは分けても分からないんです。

では、何が原因でどういうメカニズムができているのか。もちろん体の中の色々な仕組みについては、どういう原因でどういう結果というような物理学、化学の法則に従って動くのでしょうけれども、そもそもどのようにして生命のようなものが作られているかというような問いには答えられないわけですね。ことほどざるに生きているって何だろうっていう、誰しもが不思議に思う生命がある、命があるということを、近代科学的なアプローチでは答えが出せないんです。考え方ではアプローチで答えが出せないので、おそらく抜本的な発想の転換が必要だろうというわけですね。

生きているとはどういうことかという問いが持っているのは、従来の科学でアプローチできない問題点はレジュメの方にも挙げてありますが、一つは有機体が持っている、再生する、成長する、増殖するという特徴が神秘なわけですね。爪を切っても生えてきます、髪の毛を切っても生えてきます、そして新陳代謝によって私たちの体の細胞は3ヶ月もすれば全て入れ替わっているとか言われるんですけれども、その割には3ヶ月前の私と今の私は同じだっていうわけですよね。

どうやってそんなことが成立しているのか、つまり私を構成する物質的な要素を全部分解して並べてみても、構成してもできないわけですね。その中心には何だか分からないが、生命というもの、何か一体のあるそれね、これが言葉にしにくいところなんですけど、ここでは生命ですから命とか生命ということを言っておきますが、生きている生き物というのは自身というのを常に生産している。何だったらずっと成長している。

今週の水曜日に私の娘が誕生日を迎えて14歳になったんですが、昔の写真とか見ていると、これが同じ物質とは思えない。当然同じ物質、物理的には同じものではないんですよね。サイズが違いますからね。
よ。本人もこの頃の写真の記憶があるとかって言ったりするわけですから、どのようにして変わりながら同じであり続けるものを我々は維持しているのか、そしてそれを少しも不思議と思わないわけですよね。だからこの辺がもう真面目に考えると不思議で仕方がないわけです。"増殖"と書いてあるのは、もちろん細胞分裂的な事故が増殖するということもありますし、今言いましたように僕には子供がいますから、それは増殖した、自分が増殖したというのではないかもしれませんけれども、自分の情報を持ったよく似た同じようなものを作ることができる。生命というのは紡がれますから。これは近代科学の実験的にやってみようとか、分解してみようとか、何かが原因で、また親が原因で子供が生まれるみたいな雑な話だったらそれでもいいんですけど、原因と結果の、しかもその分かりやすいものでもない、数学的に厳密に表現できるものでもない、これをどう捉えたらいいのかというのが難しいわけですね。

確か九州、何処かあったかな。数年前に、高段社現代新書で、"生命は分けても分からない"という本がベストセラーになりました。10年ぐらい前だったかな。それぐらいになるかもしれませんけど、福岡陽一さんという方が書いておられる本です。軽い読み物なので読んでみてください。そこには、ハッキリとタイトルにある通り、分けても分かりません、分解しても分かりませんというようなことが書いてあります。要素、還元主義的に生命というのは研究できませんよというのは、もう誰しもが分かっているところなわけですね。黒いようですが、生きているものを分解すると、その時点で生きているものではなくなる、死んでしまうということがあるわけです。しかも、これも不思議なことに、我々生き物というのは、相手が生きているか生きていないかというのは、結構すぐに判別できるという能力も持っていますね。

何か相手が意志を持ってそうだ、心を持ってそうだ、命を持ってそうだというようなことは、我々は比較的早々に判別できるんですけれども、もちろんそれがご研修だというのが社会信じ学だったりするわけなんですけれども、そういった特徴がある生命というものを考えないわけにはいかんだろうというわけです。

生命の世界が因果的でないというのは、どういうことかといいますと、近代科学でいうところの因果の考え方というのは、AだったらB、Aが原因で原因があるから結果があるという、一方向な感じの考え方ですね。しかし、人間がすること、生命がすること、あるいはそういったものが織り出すものというのは、何が原因で何が結果になっているのかがわからない。この一つの原因が、一つの結果を及ぼすというのが、因果関係を特定するためにというような基本的なアイデアです。

因果関係というのはどうやって特定するかというと、まず原因と結果というのが確実にあって、しかも原因の方が時間的に先行している必要があります。当たり前ですね。結果が先に来て、結果があったから原因が生まれたということであれば、言葉の母義無順ですから。時間的に原因の方が先にあって、結果の方が後にあるという必要がありますし、因果関係を特定するためには、AがBになりますよというときは、Aという原因がB'にもなりません。B'にもCにもDにもなりません。AがBになりますという、すごく方向的な1対1の関係を言うわけですね。もちろんのことながら、他の要素がこっちに結びつくというのもダメです。A'もBの原因ですよ。A'もBの原因ですよ。どれかが原因ですよというようなのは、因果関係を特定したとは言いません。因果関係というのは
原因が確実にこの結果を引き起こします、という説明スタイルになっているわけです。この一方向的で時間の順序があって、一位性があるということに加えて、それらが満たされた上で、「AならばB」という考え方を持ってくるわけですね。これこそが、理論的に因果関係をつなげる発想です。

しかし、ヒュームが言ったように、「昨日も今日も朝太陽が昇ったからと言って、明日も太陽が昇る保障がどこにあるのか」。朝になれば太陽が昇るというのは、今までずっとそうだったけど、明日がそうである保障は別にないわけです。これが原因だからこれが結果、という順番で因果関係が起こるというのは、あくまで見ている人の見立てに過ぎないという考え方も、もちろんあり得るわけです。

さて、それから考えると、近代科学はわざと因果を一方向、すごく限定的な意味で使うことによって成立しています。しかし、生命や、生命をはじめとする生き生きしているもの、これらを表現するときには、近代科学的な発想というのは、いろいろとそぐわないわけですね。

生命というのは、自分も生きていますし、皆さんも生きていますし、ものすごく身近にある話なんですけれども、その神秘性は、近代科学的なアプローチではわからないということです。これを何とか理解しようというのが、今日お話しするシステム論です。

そういった意味で、生命を考えるというところからシステム論というのは始まっています。システム論という生命を考えるためのディスカッションは、いろんなところに応用できます。生きているのは必ずしも、1個人、1つの生命だけを対象にした学問ではなくて、「システム」という単位で考えましょうという考え方です。

"システム"とは、何らかのひとまとまりという意味ですが、この用語は、皆さんもよく耳にしますよね。僕も子供の頃から「システム」という単語は聞いていると思います。システムという用語がついているものをいろいろ考えてみようということで、いくつか例を挙げてみます。社会システム、情報システム、制御システムなどの仕組み的な用語がありますね。また、教育システム、経済システム、政治システムなどを言ったりすることもあります。

心理学においては、人間のシステム、行動システム、家族システムなんていう考え方もあります。一般的な社会でも生産システムや人間システムという言い方をしますし、生き物を中心に考えたら、神経システム、免疫システムなんていうのがあります。このような形を使って表現すれば、循環形や生態形というのもシステムですね。もちろん、太陽系というのもソーラーシステムです。

システムという言葉が最も身近なのは、おそらくコンピュータサイエンスのところです。オペレーティングシステム（OS）と言われますけれども、あれもシステムという言葉が付いています。データベースシステム、ネットワークシステム、コンピュータシステムもそうだし、セキュリティシステム、通信システム、といった言葉の中に入っているこの「システム」という言葉の使い方から分かることは・・・
ように、生命、生きているかのようなというのがテーマではあるんですけど、"システム"という言葉が表しているのは、何かのひとまとまりです。相互に影響し合う要素からなる全体というのがシステムの言葉の定義です。相互に影響し合う要素からなる全体。皆さん、この言葉聞いたことあるような気がしませんか？要素からなる全体のこと。

しかも、システム論を考えるときに言われるのは、生命は分けても分からないわけですから、全体は要素の層は以上です。なんとかからなる全体を表現したいんですけど、この全体は要素の層は以上であるというような定義があります。つまり、要素に分けても分からないので、要素を超えた全体ですよ。その全体的なまとまりをシステムと呼んで、そのシステムについての学をやりましょうというのがシステム論です。この要素が何でも良いというところも一つのポイントなので、押さえておいてください。

今言ったように、なんとかシステムという言葉は至るところにあるんですけれども、その言葉が示しているのは、複雑に要素が絡み合って出来上がる全体です。例えば、何でいこうかな、どのシステムでもいいんですけど、情報システムとかがあれば、インプットがある、メモリーがある、プロセスで演算処理がある、ストア、シテータ、メモリーを出しながら、ビットを置き換えて出力をするとか言ったような、一連の要素からなる処理する体系みたいなことがあったりするわけですよね。

あるいは、もっと何か、会計システムみたいなことを考えると、必要な品目だとか何だとかを人が入力すると、分かりませんけれども、予算枠のどこそこに分類されて、いろんな要素で計算して、何かお金を出したり、手続きを進めたりというようなことがあったりするわけですが、とにかく要素がたくさんあって、全体像のことを会計システムとかって読んだりするわけですよね。

このことほど作用に、一個一個の要素に分けきるのではなくて、その全体を考えましょうというのが、このシステム論、システムという言葉で表現したいことなので、これについての学を今からするんだというふうに思っておいてください。そして、これどこかで聞いたことがありますよねって今言いましたように、このシステム論というのを僕が社会シンデア部の文脈でやろうとしているのは、なぜかと言いますと、当然のことながら、集団とか社会というのもシステムだからです。

キーワードとして、クルト・レビィンの話の時にも言ったかと思うんですけれども、グループというのは、メンバーが相互に依存し合っていることが必要で、相互に依存し合うメンバーの相対、ダイナミックな全体性のことをグループとレビィンは言ったわけですね。つまり、グループというのは、生命、生きたかのような動き方をする何かなわけですね。しかもレビィンも、あるいはゲシタルトとシンデア部も言っていたように、全体は要素の相は以上という定義がありますので、ここのところをうまく心理学、社会心理学の中に取り込みたい。

そのためには、社会心理学の文脈はちょっとずれますけれども、生命を考えるというシステム論の発想と結びつけることが重要です。これについて考えてみようという話をしていきたいわけです。

この話をするにあたって、いくつか本のご案内があるんですが、話のスタートはこれです。フォン・ベルタ・ランフィーという人が書いている一般システム論というのがあります。これがシステム論がシステムという言葉で全面に出てきた最初なので、興味がある人は読んでみてください。この本の冒頭には、この一般システム論というのは学際的なものですよ。この人は生命を
スタートに考えるんですけれども、それだけではなくて、経済だとか、政治だとか、軍事だとか、社会だとか、いろんなところにこのシステムの考え方は使えますということで、一般システム論。一般というのはジェネラルな話なんですよ。何にでも応用可能な話なんですよ、ということを言っている一般システム論というのがあるわけです。

先ほども言いましたが、このシステム論という考え方そのものは物の見方です。つまり、生命をシステムとしてみましょう。グループでもいいんですが、グループをシステムとして見てみましょう。といったような見方の話になってくるので、この考え方というのは、生命現象に関わりがありそうなものの捉え方というのは、時代とともにちょっとずつ変遷していきます。

ここにもう一個別の本があるんですが、川本秀夫という人が書いているオートポイエシスという本があります。2009年ぐらいの話です。オートポイエシスという本は95年に出ています。この本の中には、今言ったシステム論の歴史的な変遷というのをまとめています。

オートポイエシスという言葉は聞き慣れないと思いますが、これもお話しするつもりです。この本の中で書いてあるシステム論の歴史的な変遷についてちょっと説明しておきますと、最初、第一世代システムというのが、今言いました一般システム理論から始まったものです。

理論としてはこうなんですけれども、どういったテーマを扱っているかというと、同的並行形です。ダイナミックイクイビリウムという言い方をしたりします。第二世代になりますと、どういうことをテーマにするかというと、これが面白いのですが、同的非平行形、第二世代システムで平行って言ったんじゃないのかと思うかもしれませんが、第二世代システムは同的非平行形です。

あるいは別名、自己組織化形というふうに呼ばれることもあります。そして、第三世代システムというのがオートポイエシスなんです。オートポイエシス。名前が難しいと思うと思うので、これはまた順応で説明しますが、これは自己創出形というふうに言われます。

このまとめ方は、川本の書いているオートポイエシスという本の中で表された考え方で、時代とともにこうやって考え方、テーマが変成してきましたよという話になっています。

ちなみに、10年くらい前に福岡新地先生の、生命は分けても分からないという本が流行ったというふうに言いましたが、あの先生が言っていらっしゃるのは同的平行形です。生き物ってのは同的平行形だよという話なんです。僕は最初に読んだときに、今更ですか?と思いました。まだ第一世代システムのことを言っているんですか?

いや、別に時代が後に来ているのが偉いとか偉くないとかっていうことは言いたいわけではないんですけど、理論というか考え方自体は進んでいっておりますので、そういったどういう変遷はあるかとか、どこに着目するか、どう変わっていったかというような視点で見ていただくとよろしいかなと思います。

さて、川本秀夫はこうやってですね、システム論の進行を時系列的に1、2、3というふうに分けました。さっきも言いましたが、これが95年くらいの話です。その後、僕が大学院生になろうとしていた頃ですから、95年、7年、8年くらいの時は言われてませんでしたが、第4世代システムの議論が出てきています。それが複雑系です。

コンプレックスシステムですね。複雑系。皆さん、この名前は聞いたことありますでしょうか?本当に2000年の頭の頃は、複雑系というタイトルの本がいっぱいあちこちにあったんです。個人的にはミシェル・ワールド・ロップの書いている複雑系という本がストーリ
ー仕立てで分かりやすいかなと思いますが、それに限らず、新書からハードモノまで、いろんなものが"複雑経"という話をやっておりました。この"複雑経"というのは、コンピュータサイエンスなども相性がいいんですけれども、いわんとしていることは、まさに生命の捉え方です。"複雑経"のテーマは、複雑なものは複雑なまま。単純化して考えようとするのがダメということと、"複雑経"という観点は、要素同士の動きは単純であったとしても、複雑な要素がエマージンとして創発されるという考え方で、表面だって見えている複雑さというのも、直接的に相互作用し合う要素同士からなる全体なんだという観点の話です。

これも、おいおいお話していきます。今日は、多分、第一世代システムの話ができませんが。

さて、4番目と5番目についてはまとめている人はいませんので、こすぎなりの考え方なんですけれども、複合システム、あるいはネットワークシステムが第5世代かなというふうに思います。ネットワークシステムというのは、ネットワークというのは基本的に要素が結びついて、ネットがワークするわけですから、結節点がつなぎ合って連動するというネットワークのシステム、その全体ということを考えましょうということなので、ネットワークシステムという話をしています。これをシステム論の視点でしゃべったのが、私の学部の時の師匠である藤沢一氏です。

そして、藤沢一氏の師匠である広田君吉というのが、一般システム論的並行形の考え方でグループダイナミックス集団心理学をやろうとした方ですね。そういう意味で、僕はずっとシステム論的な観点から社会心理学を考えるというのをずっとやってきていたわけなんですけど、流れとして、だいたいこういう感じで発展していくんだというのを先に見通しを立てておいていただければなというふうに思います。

この授業では、上から順番に解説していくつもりですが、興味がある方は、例えば川本のオートポイエシスの本を読むとかすると、ポンと先読めたりするからいいかもしれません。ちなみに、複合システムネットワークについては、藤沢一氏が"複合システムネットワーク論"という本を書いています。そこでは彼なりの別の観点からのシステム論の分類というのをやっていて、多少その辺も面白いので興味があったら読んでみてください。

ちょっと余談になりますが、第1世代システムのスタートは生物学です。第2世代システムの中心というか、考え方のポイントは、生物学、化学というところからこの考え方が出てきています。さっきも言いましたが、オートポイエシスは神経生理学の方から来た話です。こちらはコンピュータ科学です。文字が無茶苦茶ですね。最近、略して書こうとするとデジタルと言う方がやりやすいです。

この複合システムネットワーク論は、個人的には社会心理学というものですが、ネットワークシステムという意味では、これもコンピュータ科学に入っていいのかな、それとも生命科学と言った方がいいのかな? 要は、人工知能、チャットGTTだとか、生成AI、大規模言語モデルだとか、それに基づく、それの元にある機械学習、ニューラルネットワークモデルみたいなものは、このネットワークを中心にしたサイエンスですから、こういった観点は持っているはずです。

ことほどいろんなところで語られているものだから、これが社会心理学の中心ですよというのはなかなか言いにくいのですが、それぞれの観点で、どういったものの見方、何を言わんとしていたかということを、この文の見方ですから、僕ら
システムについて見ていきましょう。第1世代システムというのは、どういうものだったかというと、ホンベルタ・ランフィーが一般システム論として唱え始めたものが、この第1世代システムのアイデアです。あそこにも書いてありますように、第1世代システムの基本的なテーマは、「生きているものってダイナミックだよね」という話です。

まず、生命というものを考えたい。生命というものを考えるとき、それまでの科学的な研究というのは物理学、つまり、ものの研究をしようとすると、例えばこれはどういう仕組みで、この黒板消しがどういう仕組みでできているのかなとやったら、パーツに分解して、素性を調べて、組み合わさり方を見て、ということを切り刻んで分解するわけですね。

しかし、生命に関してはそういったアプローチができない。しかも何が違うかというと、このものというのは、研究すべきものというのは、物理学の場合は固定的なものなんですが、我々生き物みたいなことを考えると、確かに一定している状態は持っているんですけれども、ダイナミックに要素が入れ替わっていますよね。

まずこの世代システムで考えるのは、生命というのは解放系であるという発想です。解放の反対は当然のことながら閉鎖ですが、閉じた系なのか開いた系なのかということを問題にします。物の場合は物質だとか何とかのインプットとかアウトプットがありませんね。この状態のままなんですけれども、例えば生きている私を考えてみましょう。呼吸をして、酸素を取り入れて二酸化炭素を吐き出します。ご飯を食べて排泄・排尿をします。そういう物質を外から取り出してというようなことをしながら、この人間のボディを維持しているわけですね。

生命を考えるときに分けても分からないという前に、閉じた系であれば、そのままで分解しても何でもいいんだけれど、ダイナミックにインプットとアウトプットを繰り返しながら安定した状態を保っているものを研究する方法を見つけないといけないというのが、この生命を考える第一世代システムの発想なんですね。

解放システムというのは、外から物質や何かをシステムの中にインプットとアウトプットがあるということです。物質、物体、原子、電子、何でもいいですが、インプットとアウトプットを繰り返しながら、システムというのは安定していますよね。ということで、動的な、ダイナミックに開いている、その上で平行状態、安定した状態を持っているシステムを考えましょうということになっているわけです。これが、第一世代システムのものの見方です。

私はご飯を食べて何とかしてという話をしましたが、みなさん、「ホメオスタシス」（ほめようスタシス）という言葉をご存知ですかね。これは生命の維持性というものです。例えば、今、僕の子供は熱を出して寝ているんですけれども、人間が暑くなってきたら汗をかきます。寒くなってきたら、ブルブルと震えて、筋肉は震えて熱を出そうとします。そんなような感じで、生き物というのは高温動物であれば、環境の変化に応じて外界からの情報をインプットして、命令系統をアウトプットしながら、安定した体温を維持しているわけですね。

生命の基本は、開放系でインプットとアウトプットを繰り返しながら安定しているものという見方が基本なんですよということを言うわけです。それだけやったら、生命の話で終わりなんですけれども、ホンベルタ・ランフィー先生は、ここから一歩進んで、「ジェネラルシステム・セオリー」、GSTと訳されることが多いですが、一般システム理論と言うんです。この一般というのは何かというと、いろんなものにこの考え方が当ては
「う意味です。当然のことながら、GSTというのは、生命だけではなく、他のものにも当てはめて考えましょうということですね。例えば、学校の教育システムを考えると、毎年新入生が来て、卒業生をアウトプットしながら、学校という組織を維持しています。また、会計システムというのは、際による対出がありながら、会社の健全な会計状態を安定的に維持しているとも言えます。これは、ことこと作用によるものの見方の問題で、オープンなシステムで安定しているものを考えようという意味です。生命でも、学校でも、会計でも、コンピューターでも何でもいいわけです。一般的なシステム論を考え、GST、つまりGENERALということを言っています。この一般というのは、どこに強調しているポイントがあるかと言うと、オープンなシステムを数学で表現しましょうということです。数式的に表現しましょう。先ほどの現代科学の4大柱のうちの1つ、2番目に挙げました数学的な減少主義、関数主義がここにも生きています。この観点を取り入れて、ダイナミックなシステムを記述しましょうというのが一般システム論の狙いです。

ちなみに、私は先ほど大事なことを言うのを忘れていたので戻ります。生命のことを考えるとき、よくある落ち込みがちな罠がいくつかあります。そのうちの一つが全体論です。要素に還元できないのであれば、これが全体なんだと、要素に還元できない全体が何とかしているんだというだけの話にしてしまっては、ただの要素還元主義になってしまいます。全体主義になってしまってはいけません。これが一つ目の罠です。

二つ目の罠は、生命の神秘に落ち込むことです。生き物としての木の流れや、全体性の中にスピリチュアルな要素があるとか、みんなの一つにまとまろうとする気持ちがあるとか、そういう話を持っていくと、それはそれでサイエンスではなくなってしまいます。分からないけど、みんなが頑張ろうという気持ちで一つにまとまっているとか、そういう説明は何も説明したことにならないので、数学的にできれば数学的に厳密な表現を保ったまま全体の安定とは何かということを考えましょうというのが、正しいシステム論です。

ここでちょっと触れておきたいと思います。ちょっと変なことを言ったと思うかもしれませんが、これは僕も今言った単純な罠に落ちやすいところもあるので、言っておきます。ここにもう一冊本があります。それはマリコヒレオ先生という人の「複雑安定性のドグマ初めてのシステム論」という本です。これはハーデスト社が出している本で、この本は初めてのシステム論というから読むんですけど、今一体のことをすごく細かく分けています。システム論というのは、要素が複雑に組み合わさったものを考える学問です。
「て成り立つ全体」、という話をする時、深掘りしておくと全体還元主義になったり、変なスピリチュアルな傾向が出てきたりするので、そういうところに陥ってはいけないですよ。厳密に数学的な、あるいは科学的な取り組みを諦めるわけではなくて、「生きてる」っていうのは何か、ということを解きほぐそうとする気持ちが大事なのだ。そういう気持ちって言ったらおかしいね。その考え方の区分が重要なのだ、ということを強く主張しておられます。

強く主張しておられますが、個人的にはちょっと分け方が細かすぎるかな、という気もしますし。何より、この「第一世代システム」はフォン・ベルタランフィー以外にも、いろんな人、アーサー・アーケストラーだとか、いろんな人がシステム論的な観点でもって話をしているんですけれども。みんなそういう意味では、ちょっと全体観言主義だったり、ちょっと「生命の神秘ってすごいよね」みたいなのを含ませていたりということで、一概にこれが全部ダメというのもなかなか言いにくい。

言いにくいことはあるのか、分けずに語られることが多いですよ。厳密に分けてシステム論として分けるなら、今言ったような認識の違いですとか、主義主張の立て方としての「正しさ」みたいなのを論じている本がありますので、興味があったらこの本も読んでみてください。おそらくもう絶版になっていると思うんですが、図書館にはあると思います。

ちょっと横にずれたんで元戻しますけれども、「一般システム論」というのは今言ったようにダイナミックに安定するシステムというのをテーマに掲げてきました。基本的にはインプットとアウトプットがありながら安定している。そしてその話をいろんな学際的なジャンルに持っていくために数学的に表現しましょう、というのがここの目的です。

で、ベルタランフィは数学的に表現するというのをどういうふうに考えたかといいますと、伝率微分方程式を使いましょう、ということを言ったわけです。私も専門家ではないんですが、特にではないんですけれども、ベルタランフィが言おうとしたシステム論はこんなモデルです。これを書いてから言いますが、ベルタランフィはこうやってシステムを表現したらいいじゃないか、と言うんです。

この一般システム論の第三章に載っているんですけれども、「デルタT分のデルタQ1」、「デルタT分のデルタQ2」と書いてあるのは、これは微分している時間で微分している、という意味です。変化量です。何かシステムの要素の変化量というのを考えると、この要素の変化量、このQ1という要素の変化量はQ1自身を含むQ1Q2その他のこのシステムが持っている、このシステムが持っているN個の要素の関数として、Q1の変化が表現できる。このQ2の変化量というのも他の要素の関数になっている。

この要素がN個あるんだとしたら、N個の連立微分方程式でシステムというのを記述しましょう、という話なんですね。Q1Q2って何というと、これはシステムの要素です。システムの要素っていうのは何をシステムと見立てるかによって変わってきますから、それぞれのシステムにおける要素だというふうに言うしかないですね。

とにかく一般的に表現するとこうやって表現できるじゃない、ということは言いたかったんです。表現するも何も、複数の要素があってそれの関数で1個の変化、それぞれの変化量が分かるのが連立しているという話なんですけれども、これは表現はしたのはいいんですけど、言ったら表現しただけなわけです。

さっきも言ったように、このQが実態として何を表すかではなくて、考え方としてこれで行きましょうよ、という話をしているわけなんですね。ただなんでこんな話をしているかというと、ベルタランフィー曰く
論的な発想を一般化するためには、数学的に表現するしかないです。当時の数学的な限界みたいなのもあったんだと思いますが、微分しているということは、連続的な関数みたいなことを考えているわけです。それでも、連立している微分方程式の形であれば、解を見つけることができるんですよね。つまり、当時の数学では、これであったら何とか答えが見つかる。そういう形で表現するしかない、というのがベルタ・ランフィの主張だったようです。

本の最初のうちの方には、本当は二酸化炭素の数学みたいなのもやりたいとか、いろんなパターンを考えてあるんですけど、現実的に解き得るのはこれだけ、というわけです。解くも何もですね、解くっていうのはこれまた大変な話で、僕も普通の連立方程式やったら解けるんですけどね。僕は今日、午前中に本文校で行列の話をしていました。まさに連立方程式と言われてきたんですけど、この微分方程式の経過っていうのは僕も勉強が十分じゃないんですけれども。

どうやって解くかというと、「これが答えだ。Xイコール2が答えです」みたいなふうになるんではなくて、こういった表現したやつをうまく解くことができるような前提とか、過程だとか、うまい解き方ができる形に変えていくというのが、この連立微分方程式のやり方みたいです。

この関数形なら答えがはまるよ、というんで、ベル・ザ・ランフィの本の中にもですね、まずこういうものがあったとして、これが一般的な解を持つとすると、このようなテーラー展開で表現できると考えて、そういうふうなテーラー展開で表現できる連立微分方程式であれば、指数関数を含んだこういう形に書き換えることができて。

外で何が進められるかっていうと、答えは「これです」、1個の数字になるんではなくて、この式で表現できたと考えると、この係数ラムダーが0の時は安定したところに落ち着くし、0よりも大きかったらこういうパターンで動くようなシステムになっていくし、0より小さかったらこういうパターンになるようなシステムになっていくし、というような場合分けの仕方を論じるような形で、このシステムが収束していくのかとか、発散していくのかといったような一般界の出し方で表現します。

数学的な連立微分方程式の具体的な解き方みたいなことの中には踏み込みませんが、一応要素が2つだった時に解いてみると、連立方程式がこういう風に表されるから、行列式から考えてこういうパターンがあり得て、みたいなのがいくつか載ってたりします。

そして、すごく抽象的な9が何でもいいし、Nつまり数を固定してないですよ、いくらあってもいいってすごく一般的な話をしているから、なんじゃいなと。BイコールF角PEよりは具体的なのかもしれませんけど、考え方として、連立微分方程式でシステムを表現しましょうということを言ったということは認めてあげてください。

このように表現できるものとして、ベルタ・ランフィはシステムを表現したわけです。具体的なものがなかったとしても、変化量がそれぞれの要素の関数になっているよ。だから要素の組み合わさった全体なんだよというようなイメージは、この式でも十分捉えていただけるのではないかと思います。

さて、こういった形でシステムというのを表現して、この路線で生命を捉えましょうという宣言をしたのが、この一般システム論でございます。興味がある方はぜひ読んでみてください。3章のところに、システムの数学的な考察の仕方というのが書いてあります。4章、5章とかのその後の話は、こういったシステム論的な考え方をこのジャンルで応用したらどうなるかということで、心理学だとか、人間の科学と…
「システム概念だとか、生物学におけるシステム論だとか、いろんなところに応用する例なんかも書いてあります。今の数学の発展ですとか、社会侵略やろうと思ってこの本を読むと、社会侵略理論張りに中小的なことを言っているような気もしますけれども、ベルタラフィーが目指そうとしていた方向自体は、変化量で考えようとしていること、開放系としての動的な安定系を考えようとしているというようなアイディア自体は、悪くないじゃないかなと個人的に思っています。

さて、ベルタラフィーが言いました一般システム論、生物システムみたいなものを考えるときの発想なんですけれども、では、ちょっと横にずれますけど、連立微分方程式とかっていうのを知りたかったら、多分それ専門のテキストがありますが、例えば、捕食者と食べるやつと食べられるやつの動物の数がどれくらい変わるかとかっていうようなのも連立微分方程式の有名な例です。それは生態系なんです。ウォトカ・ボルテラモデルとかっていうのがあるんですけど、それも生態系なのでシステムの話になってますので、こういう表現できる世界があるんだと思ってください。

この他、このアイディアがいろんなところに展開していくわけですけれども、どういったところに応用例があるかと言いますと、一つはサイバネティクスと言われる考え方です。サイバネティクスというのは何をやってるかというと、これを考えたんですね、フィードバックループ。インプットとアウトプットするものからなるシステムというのがベルタラフィーが言ったことなんですが、このアウトプットしたものの情報をもう一度システムが取り入れて、このアウトプットした結果をもう一度取り入れて、自分でこのアウトプットの程度を調節する。これがフィードバックループを持ったサイバネティクシステムです。

名前はかっこいいんですけれども、やってることはなんというぐらいの話です。皆さんはもうすごく近代化された世界にお住まいですから、エアコンのスイッチを入れて23度とかに設定すると、あったかい風が来るでしょ。あったかい風が来ると部屋があったまってきますよね。部屋があったまってきたらエアコンは出力を弱めるはずです。それはフィードバックループがかかってるからです。

まず外の気温が寒い。その時に23度で安定させなければならないという時は、暖房を出します。暖房を出すと室温が変わります。室温が上がってきたなというと、暖房を出すのを控えるわけです。暑くなりすぎたらなんなら涼しい風を出すかもしれません。何せよ、自分のやったことをフィードバックさせることでエアコン、空調システムというのは快適な温度を保っているわけですね。

当たり前だと思うかもしれませんが、僕が中学校の頃帰っていた教室にはエアコンが付いてまして、スイッチがあるだけでした。バチンとスイッチを入れると冷房の風がドバーンと出てくるんですが、寒くなってきたら先生に「寒いです」って言ってバチンとそれを切るんです。つまりフィードバックループのシステムがなかった時代があるわけですね。

今でこそフィードバックでコントロールするのは当たり前なんですけれども、エアコンというのは冷やすものじゃない、冷房というのは冷やすものじゃない、暖房というのは温めるもんじゃないという発想からちょっと一歩進んでいるのがサイバネティクシステムというやつですね。この安定点、どこで安定してますかということの情報を、ここのアウトプットを自らコントロールしながら定めていくというシステムです。

さて、他にもこの要素全体からなる、要素の組み合わせからなる全体性という
以下のように校正いたしました。

ことで、ちょっと心理学っぽい話をしておくなら、ゲシュタルト心理学なんかも。システム論的な発想と同じところがあるわけですね。フレグナンプの法則といいますか、要素をこの形でうまい具合に置くと何か像が見えてくるとか、というのは、要素の研究ではなくて、要素と要素間の配置の関係でもって、状態が良いとか悪いとかというものが出てくるわけですから。このベルトランフィーの一般システム論の時代的には同じなんで、このデジタルルート的な発想みたいなのも今みたいなところに足を置いているわけです。

その他、システムの回想関係について考えるアイディアみたいなのも出てきます。システムというのは抽象的なようですが、何をシステムとみなすかというところなので、このシステムはどのレベルでそのシステムが成立しているかと考えるかというのはあってもいいわけですね。例えば会社組織みたいなことを考えると、会社全体がシステムです。会社の中にはきっといろんな部門があるでしょうけれども、それぞれの部門が協力し合って、生産部門、資材調達部門、営業部門、広告部門、いろんなところがあって、会社全体の安定的な利益を作っているシステムだと。

今年は何かこれが売れそうだぞというようなインプットがあったら、生産を増やすとか、生産を増やしすぎて問題が生じたんだったら、生産量を減らすとか、会社全体でバランスを取っているシステムとして会社全体を考えることもできます。だけど、会社の中にいろんな部門があって、例えば会社の中の一つの部門の生産部門があったとしたら、資材部がこれぐらい持ってきたから、それに対して作った商品をこうやって出していきましょうというようなことをするように、一つ一つの部門の小さなシステムと考えることができます。

だから、大きなシステムの中に小さなシステムがあって、例えば生産部門の中にも機械のラインを調節する部門だとか、材料を組み立てる部門があるだとか、何だとかという、いろいろな小さなサブシステムがあって、システムのネスト構造、上位システムとそれよりちょっと小さい中間のシステムと階層のシステムがあるというような、このシステムの階層構造を考えましょうというような階層システム論というのも、階層システム論というのかな、多階層関係論ですか、というのが提唱されています。

これはアーサー・ケストラーという人が、機械の中の幽霊という本の中に書いている話なんですけれども、こういった階層的なシステムの中では、中間的なシステムは、上のシステムに対しては、従属的な立場をとるんだけれども、下側を向くと、そいつが命令を発してさらにサブシステムを統合しているんだというような、それぞれのレベルでのシステム関係がありますよということを議論しています。

その発想をグループダイナミックスに持ってきたのが、広田金吉という人で、グループというのも、大きなグループ全体もあれば、グループの中のサブグループがある。サブグループの中に何だったらサブサブグループというのがある。グループというのも階層関係を成しているのだというようなことを論じています。

例えば、グループ人間関係を考えると、例えば僕は5人家族なんですが、5人家族という小さい系列という1つのシステムを持っています。その中にも親と子供は、子供は3人いますから、子供なりのコミュニケーションチャンネルがあるわけですね。親は親夫婦だけで会話するようなところもあるわけなので、小さいシステムの中には親システムと子システ
何か意見を言うとか、子システムから親システムに進行するみたいな、システムレベル同士での話があったり、その中で、個々の人間、1人1人の人間として、親としては、この冬は温泉に行きたいと言ってるけれども、僕は個人的には行った先で漫画でも読んでゴロゴロしていたりなと思ってる、個別のシステムがあったり、つまりメンバーもいろんなレベルでのシステムの関係性の中に入っているよね。これを言うようですが、一般システム論というか、同的並行形の狙いは安定ですから、それぞれのレベルで安定を目的にしていますよね。もちろん、小杉家という家族全員が幸せであるということが、幸せな状態で安定しているというのは言い訳ですけれども、その安定の中にも、夫婦関係の安定、子供関係の安定、子供の中でも姉と弟の安定とか、それぞれのレベルでの安定、先ほどの会社の例で言うともっと大きいですね。会社全体の安定もあるいは、各部門の安定もあるいは、部の中の課の関係のいろんなレベルでの安定を保ちながら、関連性を持っているものとして、そのシステムを捉えましょうという発想が必要になってくるんだよという話です。それぞれがグループダイナミックスといいますか、社会心理学の中にもシステム論的な発想というのを入れられるよと、入れて考えましょうということが重要ですよね。僕が社会心理学においてシステム論というのを持ってきているというのは、まさに今一体のところでして、個々の人間がインタラクションをします。個人の体心関係のレベルもあるし、集団レベルの話もあるし、マクロな社会のレベルでもある。それぞれがネストされた構造を持っているし、それぞれがそれぞれのレベルでの安定を求めているはずです。だから社会心理学のコアになるのは、そういった複数のシステム間の安定とか、インタラクションするところの研究が、社会心理学の本来のコアであるべきであって、インプットとアウトプットがないようなスタティックな、しかも個人の中の心理変数とか、単に外的刺激があって、それに反応する個人、その外的な刺激が社会的な刺激であったら、社会心理学とかっていうのは違うんじゃないのと思っているということです。ちなみに、第1世代システムの臨床心理学的な応用例が、今3番目に挙げてますけど、家族システム論というやつで、システム的な発想ですから、病気の治療とかっていう時も要素に還元して、家族の中で不登校になった子が出たと、上のお姉ちゃんが高校行かなくなっちゃったと、その子だけを全然ケアしても仕方がないですよ、家族関係全体の問題がその子に現れているだけなんだから、家族の中の子どもシステム、本人のシステム、子どものシステム、親子のシステム、家族のシステム、いろんなところを調整してやっていきましょうね、という方針でトリートメントしましょうという家族システム論という考え方があります。家族システム論は基本的に第1世代システムと、飛ばして第3世代システム論なんかに影響を受けた治療アプローチを取ってたりしますので、こういったシステム論の見方は、そういうところでも使えるんだということは知っておいてください。ということで、このシステム論というのはどういうことを言わんとしているかということの方向性はお分かりいただけたかなと思うんですが、第1世代システム論の限界とか、次にどういったことが問題になってきたかということをちょっと指摘しておきたいと思います。まずは、この連立微分方程式で書いたことによる限界です。何かと言いますと、システムはオ
システムがダイナミックに関わって安定していきますよ、というところまではよかったんですが、ここです。サイズNの連立微分方程式になっていますね。この要素の数を限定してしまうというのが、ベルザ・ランフィードの名字的には語られないのですが、限界としてあります。皆さんはこの辺について、そのようにお考えかなと思うんですが。

最近生成AIってすごいじゃないですか。生成AI、僕もよくチャットGPTだとか、クローダーとかと会話しながら、あるいはこの技術力、講義力を取るのも彼ら、それらの力を借りながらやるんですけど、すごくそれっぽい会話ができてます。まるで生きた人間かのように、知能があるかのように、生命かのようにやりとりができるんですけれども、彼らは学習したセット、膨大な学習したセットがあって、その中の組み合わせから言葉を作って出力しています。

膨大な学習セットがありますから、例えば僕はこんなプログラム書きたいんだとかって言ったら、すでに誰かが書いたことがあるプログラム、効率の良さそうなプログラムですとかをうまく組み合わせて、そしてこういう問いにはこういう答え方をしたらよいだろうというパターンから考えてそれっぽい答えを持ってきます。

しかし、あくまでもそこはスタティックな、もう学習の終わった、膨大ではありますけど、固定的なデータセットっていうのはあるわけです。これは人間との本質的な違いじゃないかなと、僕は個人的には考えています。人間というのは、原理的には、このインプットするNのサイズが、ここがダイナミックに変わるんですよね、きっと。僕の考え方なのでそうではないと言われるかもしれませんが。

人間は新しい知識、新しい言葉を作ろうということもできるし、自分のボキャブラリを増やすために本を読むこともできる、情報を仕入れることができる、データセットをダイナミックに毎回変えていくことができるわけです。いやいや、もう最近の子は本は読まないよとか、インプットせずに自分の小さな凝り固まった知識の中だけで議論するやつばっかりですよ。

人間がどれだけ頑張ったとしても、チャットGBTが学習したデータセットほどもテキストは読めませんよということはあるかもしれませんが、原理的にはそこが可変であるということが、僕は人間の大事なポイントなんだろうなと、チャットGBTなんかと違うところなんだろうなと思います。

生きている、生命を持っているということは、これが変わり得るということなんじゃないかなと思うんです。グループとかということを考えたとしても、例えばそうですね、野球のチームっていうのは9人で野球しますから、メンバーは9って決まってるんです。確かに9っていう、この9と引っ掛けましてね、9人でやりますというのは決まっていますが、本来有機的なグループっていうのは、メンバーが出たり入ったりしたとしても安定的な何かをやっているはず、やっているもののことを言うんじゃないかなと思うわけですね。

ベルタ・ランフィーが連立微分方程式の形で変えたときに、要素のサイズを固定しちゃった、固定せざるを得なかったということは、一つの大事なシステムが考えるべき問題の限界だったかなと思います。もう一つは先ほども言いましたが、連立微分方程式を選んだこと、そのものですね、本当に連続的なものなのかを理解するということは、積分したら元に戻るみたいな意味もありますからね。つまりニュートンの世界の微分、積分っていうのは、逆の運動がそこに成立しているわけです。ということは、時間を微分してっていうことができるということは、積分したら元の状態に戻るみたいな対象
性があるわけなんですけれども、さて、システムあるいは生命が考えなければいけない「時間の方向性」っていうのは、そんな対象性ってありましたっけ？っていう話ですね。ちなみに、今のこの時間の非対象性の話については、第二世代システム化学のところでエントロピーの増大に関する問題で、方向が一方向であるというようなテーマにつながってきます。

サイズNが固定されているということとも関係しますが、ベルタランフィーのシステムの限界の一つは、増殖・修復・創出に言及できなかったことですね。生命とは何かという話をしたときに、僕は例として、「再生するよ。爪を切っても髪の毛を切っても伸びてくるよ」とか、「自分のコピーを作る生き物というのは自分のコピーを作ることができますよ」というような話をしましたが、ベルタランフィーのこのシステムでは、そういったことまでは言及できていません。

あくまでもオープンで安定したシステムを数式的に書くところまでなので、これはまだちょっと生命を語るという意味では、限界があるわけです。先ほどもちょっと言いましたけど、正規論・勇気対論、全体還元主義に陥らないように注意も必要です。こういったシステム論で、全体というのは要素の層は異常である、非常に聞こえはいいんですが、それを言ってしまうと、「だから全体なのだ」というところに陥ってしまいがちなので、しかもその全体になろうとしているのは、そういう方向に行こうとしたからだというような目的性だとか、それこそ最初に否定されるべきだった、一直線の因果性を含めてしまいがちになるので、そこは注意した方がいいでしょう。

こういった限界はあるんですけれども、ベルト・ランフィーが一般システム論を立てたおかげで、一つのものの見方が成立したというのは評価すべきところではないかなと思います。皆さんはこの第1世代システム、同的平行形というのをどのように捉えていただけましたでしょうか。自分なりにオープンなシステム、オープンで安定しているシステムというのを考えてみるというのが良いかなと思います。ちょっと考えてみてください。

この後につながる話でもあるので、第2世代の話もちょっと入る話で注意が必要なんですが、このシステム論の考えるべき問題、ベルト・ランフィーが考えてなかった問題の限界がもう1つありまして、それは何かというと観察者問題です。今、システムというのはものの見方ですと、何かをシステムとして捉えると、「これがインプットで、これがアウトプットだ」と言いました。ベルト・ランフィーはサイズNの関数を考えて分析しようとしていますよという話もしました。僕はこのサイズがNが変わってもいいんじゃないですかねという話をしましたが、そういった話の中にひっそりと隠れているのは、何を全体と見なすかということがテーマになっているんですが、その時同時に、「このシステム、これがシステムであると見なした人はどこにいますか？」という問題があります。

空調のシステムとかと言いますと、エアコンを取り出して外気をインプットして、設定した温度に合うようなアウトプットをやりますと言った時に、「この空調がシステムですよ」と言っている人は、明らかにそのシステムの外にいているわけですね。システムの外にいながら「23度です」という安定点を外から設定しているわけです。そうなると、「そいつ誰やねん」という話になります。いやいや、エアコンなんだからあなたでしょと言われるかもしれませんが、確かにエアコンだったら僕が設定して暖かい寒いというふうにしたらそれでいいんですけれども、生き物のことを考えてみてください。皆さん生きてますよね。左の一人も生きてますね。
皆さんの安定というのは、誰が決めたんですか？僕があなたの体温は23度ですよ、と決めたら、23度は低すぎますけどね。あなたの体温、これぐらいですよ、と決めたわけではないですね。あなたの幸せは、この辺ですよ、と決めたわけでもないですね。

自分というのが、このボディの物理的なものだけだと皆さんはお考えでしょうか？もしそうだとしたら、自分の境界線ははっきりしていますが、自分、事故、生命というのは、少なくとも物理的なものだけじゃないですよね。この辺はうまく説明しないといけないんですが、いつしかここに小さいおじさんがいると考えてみたいな話をしましたけど、社会信念を考えたらいいのか。

実態というのは、社会的に実在していればいいのであって、物理的な実態が必要ではない。グループというのは目に見えない全体なので、見える必要はない。この枠の中に入った人がグループという定義でも、言っちゃいいけれども、そんなはずないじゃないですか。物理的なグループの境界というのは、物理的に決められるものではないはずです。

ということは同時に、人間の私の境界というのも、別にこの物理的な基盤は持っていて、なるべく皆さんの像がフィックスしやすいような形はありますが、形がなかったとしても、その人個人のシステムというのは、実在していればいいわけなので、ここに固定されているわけではないわけです。

何が言いたいかというと、生命とは、生きているとは、グループであるとは、自分があるというのは、外から観察されているわけじゃないんですよ。自分で決めているんです。自分の安定点というのも自分で作っているんです。この事故を持っているというのが、この生命の特別な性質なので、この事故をどう扱いますかというのが、第2、第3世代のシステムに関わってきます。

そのつもりで、次回以降のお話を聞いていただければと思います。最後に、申し訳程度に、社会心理学における動的並行形という話題を用意しておきました。

一つは、家族システムの話は今も言った通りです。臨床経営の家族システム論というものの中では、このベルタランフィシステムの観点で、因果的な直線的な要素、還元主義的な治療、介入をしないでおこうというような方針として運用されています。治療的な話がなかったとしても、社会心理学におけるグループや、家族の考え方もあるでしょうから、社会心理学における動的な並行形という意味で、家族やグループというのを考えることもできるでしょう。

それともう一つは、いつしか社会心理学の歴史の中でお話ししましたが、ハイダーやニューカム、オズグット・トゥカネンバーム・フェスティンガーなどが考えた認知均衡理論というのがありましたね。バランス理論です。インバランスしていたらバランスした状態に戻そうとしますよ。バランス理論というのも、社会心理学における動的並行形の観点を持っているものですね。

いろんな認知的な要素があるんだけど、それが安定的な全体性になるように、情報のインプット・アウトプットをしながら、安定的な心の状態を保とうとする考え方なわけです。なので、バランス理論というのも、第8世代システム的な発想を持った考え方なんじゃないかなと思います。

さて、もう一つは多変量解析ですね。社会心理学の具体的な研究方法として、社会調査を行って、つまりアンケート調査、調査表によるデータを対応にとって、それを分析して、何かモデルを立てるというような研究が主流です。このことに関する批判は、また後半の方でやっていこうかなと思うんですけれども、多変量解析、社会調査による分析方法というのは、アンケート調査表というデータ
「に含まれる安定的な構造を取り出そう」という考え方です。その取り出した構造によって、このアンケート調査の空間の中には、「こういう要素、こういう要素、そしてこういう要素がありましたよ」というような言い方をします。しかし、「あの因子がある」というのは、具体的な厳密な話をすると、一人一人の心の中に因子があるわけではありません。社会だとか言語空間の中に、因子という安定的な構造が存在しています。そこに調査協力者、あるいは被験者、といった人を作用させると、人に項目のセットを与えると、因子があるような形に変換されます。

数学でいうところの「写像の話」だとか、「行列による座標変換の話」といったものです。言語空間における安定的な要素間関係を記述しましょう、というのが大変重要なのです。主成分分析とか因子分析というのは、言語空間における要素間関係を記述するための方法です。なので、こういった調査方法なんかも基本的には「動的平行形」です。言語空間というのはきっとこういう因子で安定しているでしょう。

そうすると、人が社会問題に対して意識を持ったら、「こういう3つの要素に分かれます」といったような形でモデル化されているのが非常に理解しやすいです。そういったものとして、社会学の中でも「動的平行形」の考え方が生きている、ということについて、最後にちょっと触れておきたいと思います。

はい、今日は「システム論の導入」ということで、第1世代システムをしっかりと学んだつもりです。次回は第2、第3世代のシステムについてお話を進めていきたいと思います。その時のテーマは先ほども言いましたが、「自己の境界をどうするか」というようなことと、それに含めて「観察者の問題」なんかが入ってきますので、どうぞお楽しみに。はい、それでは今日の話はここまでにしておきましょう。お疲れ様でした。

\printbibliography[title=引用文献]

\end{document}

